\chapter[Implementation Governance, Change Management, dan Requirements]{Implementation
Governance, Architecture Change Management, dan Requirements Management}
\label{ch:governance-change}
\setlength{\emergencystretch}{2em}

\section*{Tujuan Pembelajaran}
\begin{itemize}
  \item Menjelaskan \textbf{Implementation Governance} (\emph{tata kelola
        implementasi}, TOGAF Fase G): \emph{Architecture Contract},
        \emph{Compliance Review}, dan \emph{Dispensation}.
  \item Menjelaskan \textbf{Architecture Change Management} (\emph{manajemen
        perubahan arsitektur}, TOGAF Fase H): \emph{change driver} dan klasifikasi
        perubahan (\emph{Simplification}, \emph{Incremental}, \emph{Re-architecting}).
  \item Menjelaskan \textbf{Requirements Management} sebagai \emph{pusat} ADM:
        \emph{Requirements Repository}, \emph{traceability}, dan aliran perubahan
        antar-fase.
  \item Menutup siklus ADM sehingga arsitektur menjadi praktik yang \emph{hidup
        dan terus dikelola}---bukan dokumen sekali jadi---serta mengaitkannya
        dengan tata kelola dan kepatuhan di Indonesia.
\end{itemize}

Tujuan pertama dan kedua mendukung CLO 6; tujuan ketiga mendukung CLO 2 dan
CLO 3; tujuan keempat mendukung CLO 6 dan CLO 7 karena mahasiswa merumuskan tata
kelola, risiko, dan implikasi manajerial yang dapat dipertahankan secara akademik.

\section{Ringkasan Awal Bab}

Bab~\ref{ch:implementation-migration} sudah menjawab \emph{apa} arsitektur
targetnya dan \emph{bagaimana} membangunnya melalui peta jalan. Namun, rencana
sebaik apa pun bisa melenceng ketika dieksekusi, dan keadaan terus berubah setelah
sistem beroperasi. Bab ini menjawab dua pertanyaan penutup: \emph{bagaimana
memastikan pembangunan tetap setia pada rancangan} (\emph{governance}) dan
\emph{bagaimana arsitektur menyesuaikan diri ketika keadaan berubah}
(\emph{change management}). Keduanya diikat oleh satu benang yang selalu berada di
pusat ADM: \textbf{Requirements Management}.

Tiga gagasan intinya: \textbf{Fase G (Implementation Governance)} memastikan proyek
implementasi patuh pada arsitektur target melalui \emph{Architecture Contract} dan
\emph{Compliance Review}; \textbf{Fase H (Architecture Change Management)} memantau
\emph{change driver} lalu memutuskan apakah cukup ditangani lewat governance atau
harus memicu \emph{siklus ADM baru}; dan \textbf{Requirements Management} menjaga
\emph{traceability} setiap requirement dari motivasi sampai solusi yang berjalan.
Karena bab ini menutup rangkaian metode, keluarannya---\emph{architecture charter},
matriks traceability, dan protokol perubahan---mengisi bagian \textbf{Tata Kelola},
\textbf{Risiko}, dan \textbf{Implikasi Manajerial} pada naskah.

\section{Pendahuluan}

Banyak organisasi memperlakukan arsitektur sebagai proyek: dokumen disusun,
dipresentasikan, lalu dilupakan begitu pembangunan dimulai. Akibatnya bisa
ditebak: tim implementasi mengambil jalan pintas, sistem yang terwujud berbeda
dari yang dirancang, dan setahun kemudian tak seorang pun ingat mengapa suatu
keputusan diambil. Di sinilah \emph{governance} dan \emph{change management}
menutup celah itu. Fase G membuat kepatuhan arsitektur \emph{terukur} sepanjang
pembangunan; Fase H membuat arsitektur \emph{tangguh menghadapi perubahan}; dan
Requirements Management membuat seluruh rantai keputusan \emph{dapat ditelusuri}.

Bagi makalah, ketiga hal ini mengubah rekomendasi dari sekadar rancangan menjadi
\textbf{praktik yang dapat dipertanggungjawabkan}. Gagasan \emph{architecture
governance}, \emph{architecture contract}, dan pemantauan perubahan mengikuti
TOGAF ADM Fase G, H, dan Requirements Management \cite{togaf10}. Perspektif bahwa
tata kelola dan fondasi eksekusi adalah kapabilitas organisasi yang dibangun
sengaja---bukan formalitas administratif---selaras dengan gagasan \emph{operating
model} dan \emph{IT governance} \cite{ross2006enterprise,ross2019designed}. Praktik
pengiriman perangkat lunak yang berkelanjutan (\emph{continuous delivery}) pun
menuntut agar perubahan dikelola melalui umpan balik, bukan dihindari
\cite{newman2021building}.

\section{Skenario Pembuka dan Pertanyaan Pemandu}

Peta jalan rel transfer Bank Wanua sudah disetujui dan pembangunan dimulai. Kini
muncul ketegangan yang khas: tim pengembang ingin bergerak cepat menuju
\emph{go-live} SNAP BI, sementara arsitek khawatir gateway yang dibangun mulai
menyimpang dari rancangan keamanan. Di tengah pembangunan, Bank Indonesia
mengumumkan pembaruan spesifikasi SNAP BI, dan OJK meminta bukti \emph{traceability}
kontrol AML dari kebijakan sampai kode. Pertanyaannya bukan lagi ``apa yang
dibangun'', melainkan ``bagaimana memastikan yang dibangun tetap benar dan tetap
relevan''.

\subsection*{Pertanyaan Pemandu}
\begin{enumerate}
  \item Bagaimana memastikan proyek implementasi patuh pada arsitektur target?
  \item Apa isi \emph{Architecture Contract}, dan kapan sebuah penyimpangan boleh
        diberi \emph{Dispensation}?
  \item Ketika muncul perubahan (regulasi, teknologi, kebutuhan baru), bagaimana
        menilai apakah cukup ditangani lewat governance atau perlu siklus ADM
        baru?
  \item Bagaimana setiap requirement ditelusuri dari pemangku kepentingan sampai
        solusi yang berjalan?
  \item Tata kelola, metrik, dan risiko apa yang membuat arsitektur tetap hidup
        dan patuh?
\end{enumerate}

\section{Mengapa Bab Ini Penting bagi Pembaca Non-TI}

Justru pembaca non-TI---direksi, komite audit, regulator, pemilik risiko---lah yang
paling berkepentingan pada bab ini, sebab di sinilah arsitektur berjumpa dengan
\emph{akuntabilitas}. Tanpa istilah teknis, seorang analis dapat menjelaskan
\emph{siapa bertanggung jawab atas kepatuhan} (\emph{governance}), \emph{bagaimana
keputusan perubahan diambil} (\emph{change management}), dan \emph{bagaimana setiap
kapabilitas dapat dilacak kembali ke alasannya} (\emph{Requirements Management}).
\emph{Traceability} inilah yang diminta auditor dan regulator: bukti bahwa sistem
yang berjalan benar-benar mewujudkan kebijakan dan kebutuhan yang dijanjikan.

\section{Konsep Inti dan Glosarium Mini}

\subsection{Fase G: Implementation Governance}

\textbf{Fase G} memastikan proyek implementasi sesuai dengan arsitektur target.
Instrumen utamanya adalah \emph{Architecture Contract} (kontrak arsitektur)---
kesepakatan formal antara fungsi arsitektur dan tim implementasi yang menetapkan
\emph{apa arti patuh}: standar, prinsip, dan batasan yang harus dipenuhi. Sepanjang
pembangunan dijalankan \emph{Architecture Compliance Review} (tinjauan kepatuhan):
apakah yang dibangun sesuai rancangan? Bila ada penyimpangan, ada dua jalan---
\emph{rework} (perbaikan) agar kembali patuh, atau \emph{Dispensation} (dispensasi):
pengecualian berbatas waktu yang disetujui \emph{Architecture Board} (dewan
arsitektur) dengan alasan dan rencana penyelesaian. Fase G-lah yang menjaga agar
peta jalan Fase E/F benar-benar dieksekusi seperti dirancang.

\subsection{Fase H: Architecture Change Management}

\textbf{Fase H} berlaku setelah solusi beroperasi. Ia memantau \emph{change driver}
(pemicu perubahan)---teknologi baru, kebutuhan bisnis baru, regulasi baru, masalah
kinerja---lalu menilai dampaknya dan \emph{mengklasifikasi} perubahan menjadi tiga
jenis \cite{togaf10}: (1) \emph{Simplification} (penyederhanaan): mengurangi
biaya/kompleksitas, ditangani lewat change management biasa; (2) \emph{Incremental}
(inkremental): memperluas atau menambah nilai dalam arsitektur yang ada, mungkin
menuntut sebagian ADM; dan (3) \emph{Re-architecting} (penataan ulang): perubahan
mendasar yang \emph{memicu siklus ADM baru}, kembali ke Fase A. Keputusan inti
Fase H: apakah perubahan cukup dikelola dalam governance, atau harus memulai
iterasi arsitektur baru---sekaligus menilai \emph{value realization} (realisasi
nilai), yakni apakah arsitektur yang berjalan benar-benar memberi kapabilitas yang
dijanjikan.

\subsection{Requirements Management: Pusat ADM}

\textbf{Requirements Management} bukan sebuah fase, melainkan \emph{pusat} roda ADM
yang terhubung ke \emph{semua} fase (Gambar~\ref{fig:adm-ghr}). Setiap fase
\emph{menghasilkan} dan \emph{memakai} requirement; perubahan pada requirement
mengalir ke fase-fase terkait. Prosesnya menyimpan requirement dalam
\emph{Architecture Requirements Repository} (repositori requirement),
memprioritaskannya, dan---
yang terpenting---menjaga \emph{traceability} (keterlacakan): setiap requirement
tertaut ke elemen arsitektur dan work package yang memenuhinya, dan kembali ke
motivasi yang melatarinya (Bab~\ref{ch:strategi-bisnis}). \emph{Traceability}
inilah yang membuat arsitektur \emph{dapat diaudit dan dipertahankan}.

\begin{figure}[htbp]
\centering
\scalebox{0.92}{
\begin{tikzpicture}[
  font=\footnotesize\bfseries,
  phase/.style={circle, draw=praditagreen, line width=0.3mm,
    fill=praditagreen!12, minimum size=10mm, inner sep=0pt},
  hi/.style={circle, draw=blue!70!black, line width=0.6mm, fill=blue!35,
    minimum size=10mm, inner sep=0pt},
  reqman/.style={circle, draw=orange!80!black, line width=0.5mm, fill=orange!22,
    minimum size=20mm, inner sep=1pt, align=center, font=\scriptsize\bfseries},
  prelim/.style={circle, draw=gray!70, line width=0.3mm, fill=gray!15,
    minimum size=11mm, inner sep=0pt, align=center, font=\scriptsize\bfseries},
  ring/.style={-Latex, draw=praditagreen!70, line width=0.3mm},
  spoke/.style={draw=orange!60, line width=0.25mm},
]
  \node[reqman] (rm) at (0,0) {Requirements\\Management};
  \foreach \ang/\lab in {90/A, 45/B, 0/C, -45/D, -90/E, -135/F}
    \node[phase] (\lab) at (\ang:2.4cm) {\lab};
  \node[hi] (G) at (180:2.4cm) {G};
  \node[hi] (H) at (135:2.4cm) {H};
  \node[prelim] (P) at (90:3.8cm) {Prelim.};
  \foreach \lab in {A,B,C,D,E,F,G,H} \draw[spoke] (\lab) -- (rm);
  \foreach \a/\b in {A/B, B/C, C/D, D/E, E/F, F/G, G/H, H/A}
    \draw[ring] (\a) to[bend left=12] (\b);
  \draw[ring, draw=gray!70] (P) -- (A);
\end{tikzpicture}
}
\caption{Fokus bab ini pada roda TOGAF ADM: \textbf{Fase G} (Implementation
Governance) dan \textbf{Fase H} (Architecture Change Management) disorot, dengan
\textbf{Requirements Management} sebagai pusat yang terhubung (jari-jari) ke semua
fase.}
\label{fig:adm-ghr}
\end{figure}

\subsection*{Glosarium Mini}
\begin{description}
  \item[Architecture Contract] Kesepakatan formal yang menetapkan arti kepatuhan
    bagi tim implementasi.
  \item[Compliance Review] Pemeriksaan apakah yang dibangun sesuai arsitektur
    target.
  \item[Dispensation] Pengecualian berbatas waktu atas penyimpangan, dengan alasan
    dan rencana penyelesaian.
  \item[Change driver] Kejadian (teknologi, bisnis, regulasi) yang menuntut
    penilaian ulang arsitektur.
  \item[Klasifikasi perubahan] \emph{Simplification}, \emph{Incremental}, atau
    \emph{Re-architecting}.
  \item[Architecture Requirements Repository] Tempat requirement disimpan,
    diprioritaskan, dan ditelusuri.
  \item[Traceability] Tautan requirement $\rightarrow$ arsitektur $\rightarrow$
    solusi $\rightarrow$ motivasi.
\end{description}

\section{Kerangka Berpikir}

Alur berpikir bab ini melingkar, bukan lurus (Gambar~\ref{fig:gov-flow}).
Arsitektur target dan peta jalan masuk sebagai \emph{Architecture Contract};
governance menguji kepatuhan pembangunan; solusi yang patuh dirilis; setelah
beroperasi, perubahan dipantau dan diklasifikasi; perubahan besar memicu siklus
ADM baru, sedangkan perubahan kecil kembali ke governance. Requirements Management
berada di tengah, menjaga \emph{traceability} di setiap langkah.

\begin{figure}[htbp]
\centering
\scalebox{0.84}{
\begin{tikzpicture}[
  box/.style={rectangle, rounded corners=1.5mm, draw=blue!60!black,
    fill=blue!8, thick, align=center, text width=27mm, minimum height=11mm},
  hub/.style={rectangle, rounded corners=1.5mm, draw=orange!80!black,
    fill=orange!16, thick, align=center, text width=27mm, minimum height=11mm},
  link/.style={-Latex, thick, draw=blue!60!black}
]
  \node[box] (con) {Architecture Contract\\(dari peta jalan)};
  \node[box, right=7mm of con] (gov) {Compliance Review\\(Fase G)};
  \node[box, right=7mm of gov] (dep) {Solusi patuh\\dirilis};
  \node[box, below=8mm of dep] (mon) {Pantau change driver\\(Fase H)};
  \node[box, left=7mm of mon] (cls) {Klasifikasi\\perubahan};
  \node[hub, left=7mm of cls] (req) {Requirements\\Management (pusat)};
  \draw[link] (con) -- (gov); \draw[link] (gov) -- (dep); \draw[link] (dep) -- (mon);
  \draw[link] (mon) -- (cls); \draw[link] (cls) -- (req);
  \draw[link, dashed] (req) to[out=90,in=180] node[above,font=\scriptsize]{ADM baru / penyesuaian} (con);
\end{tikzpicture}
}
\caption{Alur berpikir governance \& change management untuk makalah: siklus
tertutup dari contract, compliance, rilis, pemantauan, klasifikasi, hingga kembali
ke Requirements Management di pusat.}
\label{fig:gov-flow}
\end{figure}

\section{Fase G: Contract, Compliance Review, dan Dispensation}

Implementation Governance bekerja sebagai lingkar umpan balik
(Gambar~\ref{fig:governance-loop}). Rancangan arsitektur menurunkan sebuah
\emph{Architecture Contract} yang mengikat proyek implementasi. Selama proyek
berjalan, dijalankan \emph{Compliance Review}. Bila hasilnya \emph{patuh}, solusi
lanjut ke rilis; bila \emph{menyimpang}, penyimpangan dinilai: kalau dapat
dibenarkan sementara, Architecture Board memberi \emph{Dispensation} berbatas
waktu; kalau tidak, proyek melakukan \emph{rework} dan ditinjau ulang. Lingkar ini
menjaga agar kualitas arsitektur tidak tergerus tekanan jadwal.

\begin{figure}[htbp]
\centering
\scalebox{0.95}{
\begin{tikzpicture}[
  font=\footnotesize,
  gov/.style={rectangle, rounded corners=1.5mm, draw=blue!60!black, fill=blue!10,
    thick, align=center, text width=28mm, minimum height=10mm},
  dec/.style={diamond, aspect=2, draw=blue!60!black, fill=blue!6, thick,
    align=center, inner sep=1pt, text width=20mm},
  rels/.style={rectangle, rounded corners=1.5mm, draw=praditagreen, fill=praditagreen!12,
    thick, align=center, text width=24mm, minimum height=9mm},
  flow/.style={-Latex, thick, draw=gray!70}
]
  \node[gov] (con) {Architecture Contract};
  \node[gov, right=14mm of con] (rev) {Compliance Review};
  \node[dec, right=14mm of rev] (d) {Patuh?};
  \node[rels, above right=6mm and 8mm of d] (rel) {Rilis solusi};
  \node[gov, below right=6mm and 8mm of d] (disp) {Dispensation / Rework};
  \draw[flow] (con) -- (rev);
  \draw[flow] (rev) -- (d);
  \draw[flow] (d) -- node[above,font=\scriptsize]{ya} (rel);
  \draw[flow] (d) -- node[below,font=\scriptsize]{tidak} (disp);
  \draw[flow] (disp.west) to[out=180,in=-90] node[left,font=\scriptsize]{tinjau ulang} (rev.south);
\end{tikzpicture}
}
\caption{Lingkar Implementation Governance (Fase G): Architecture Contract
$\rightarrow$ Compliance Review $\rightarrow$ keputusan patuh; penyimpangan
ditangani lewat Dispensation berbatas waktu atau rework lalu ditinjau ulang.}
\label{fig:governance-loop}
\end{figure}

\section{Fase H: Memantau dan Mengklasifikasi Perubahan}

Begitu solusi beroperasi, arsitektur langsung berhadapan dengan dunia yang terus
berubah. Fase H mengubah perubahan dari ancaman menjadi proses yang terkelola
(Gambar~\ref{fig:change-mgmt}). Sebuah \emph{change driver} dinilai dampaknya, lalu
diklasifikasi. \emph{Simplification} dan sebagian besar \emph{Incremental} cukup
ditangani lewat change management---memperbarui contract dan menjalankan Compliance
Review ulang. \emph{Re-architecting} yang mendasar memicu \textbf{siklus ADM baru}:
kembali ke Fase A untuk merumuskan visi arsitektur berikutnya. Aturan praktis The
Open Group membantu memilah: bila sebuah perubahan menyentuh \emph{dua pemangku
kepentingan atau lebih}, ia cenderung menuntut penataan ulang dan masuk kembali ke
ADM; bila hanya satu, ia lebih cocok ditangani lewat change management
\cite{togaf10}. Dengan begitu, arsitektur tidak pernah usang tanpa disadari; setiap
perubahan tercatat, terklasifikasi, dan tertelusur.

\begin{figure}[htbp]
\centering
\scalebox{0.9}{
\begin{tikzpicture}[
  font=\footnotesize,
  ev/.style={rectangle, rounded corners=1.5mm, draw=purple!60!black, fill=purple!8,
    thick, align=center, text width=26mm, minimum height=10mm},
  cls/.style={rectangle, rounded corners=1.5mm, draw=blue!60!black, fill=blue!8,
    thick, align=center, text width=26mm, minimum height=9mm},
  act/.style={rectangle, rounded corners=1.5mm, draw=praditagreen, fill=praditagreen!12,
    thick, align=center, text width=30mm, minimum height=9mm},
  flow/.style={-Latex, thick, draw=gray!70}
]
  \node[ev] (drv) {Change driver\\(teknologi, bisnis, regulasi)};
  \node[cls, right=12mm of drv] (imp) {Nilai dampak \&\\klasifikasikan};
  \node[act, above right=8mm and 12mm of imp] (small) {Simplification / Incremental\\$\rightarrow$ change management};
  \node[act, below right=8mm and 12mm of imp] (big) {Re-architecting\\$\rightarrow$ siklus ADM baru (Fase A)};
  \draw[flow] (drv) -- (imp);
  \draw[flow] (imp) -- (small);
  \draw[flow] (imp) -- (big);
\end{tikzpicture}
}
\caption{Architecture Change Management (Fase H): change driver dinilai dan
diklasifikasi; perubahan kecil ditangani lewat change management, perubahan
mendasar memicu siklus ADM baru.}
\label{fig:change-mgmt}
\end{figure}

\section{Input, Langkah, dan Output Tiap Fase}

Agar tiga fase penutup ini dapat \emph{diterapkan}, bukan sekadar dipahami, tiap
fase perlu diperjelas \emph{input} (apa yang masuk), \emph{langkah} (apa yang
dikerjakan), dan \emph{output} (apa yang dihasilkan). Tabel~\ref{tab:io-fase-g}--
\ref{tab:io-rm} merangkum input, langkah, dan output kanonik tiap fase, lalu
melengkapinya dengan dua contoh konkret---restoran pemesanan online dan rel
transfer Bank Wanua---sehingga polanya terlihat sama meski domainnya berbeda.

\begin{table}[htbp]
\centering
\scriptsize
\caption{Input, langkah, dan output Fase G (Implementation Governance).}
\label{tab:io-fase-g}
\begin{tabularx}{\textwidth}{@{}p{0.17\textwidth} X@{}}
\hline
\textbf{Aspek} & \textbf{Uraian} \\
\hline
Input & \emph{Architecture Contract}, \emph{Architecture Definition Document},
\emph{Architecture Roadmap}, dan \emph{Implementation \& Migration Plan} (keluaran
Fase E/F). \\
Langkah & (1) Konfirmasi lingkup \& prioritas deployment; (2) siapkan sumber daya
\& keterampilan; (3) bimbing pembangunan agar sesuai arsitektur; (4) jalankan
\emph{Architecture Compliance Review}; (5) go-live, lalu \emph{post-implementation
review}. \\
Output & \emph{Architecture Contract} yang ditandatangani, \emph{Compliance
Assessment}, \emph{Change Request}, solusi patuh yang ter-deploy, dan
\emph{Architecture Repository} yang terisi. \\
Contoh (Restoran Go-Online) & \emph{Input:} Architecture Contract aplikasi pemesanan
(standar keamanan pembayaran, integrasi POS \& dapur). \emph{Langkah:} bimbing
\emph{build} integrasi $\rightarrow$ Compliance Review pra-peluncuran $\rightarrow$
luncurkan lalu review. \emph{Output:} Compliance Assessment ``lulus bersyarat'',
Dispensation modul pembayaran 1 bulan, Change Request integrasi internal, aplikasi
live. \\
Contoh (Bank Wanua) & \emph{Input:} Architecture Contract rel transfer (standar
keamanan gateway, batas integrasi core). \emph{Langkah:} konfirmasi lingkup
$\rightarrow$ bimbing \emph{build} gateway/adaptor $\rightarrow$ Compliance Review
WP3 pra go-live $\rightarrow$ go-live SNAP BI lalu review. \emph{Output:} Compliance
Assessment WP3, Dispensation adaptor core 3 bulan, Change Request lapisan audit, rel
transfer live. \\
\hline
\end{tabularx}
\end{table}

\begin{table}[htbp]
\centering
\scriptsize
\caption{Input, langkah, dan output Fase H (Architecture Change Management).}
\label{tab:io-fase-h}
\begin{tabularx}{\textwidth}{@{}p{0.17\textwidth} X@{}}
\hline
\textbf{Aspek} & \textbf{Uraian} \\
\hline
Input & \emph{Change Request} (dari teknologi, bisnis, atau regulasi),
\emph{Architecture Contract}, \emph{Compliance Assessment}, serta arsitektur \&
requirement yang sedang berjalan. \\
Langkah & (1) Pasang pemantauan \& proses \emph{value realization}; (2) tangkap
\emph{change driver}; (3) nilai dampak \& kelola risiko; (4) klasifikasikan
perubahan (\emph{Simplification}, \emph{Incremental}, atau \emph{Re-architecting});
(5) tangani lewat governance atau picu siklus ADM baru. \\
Output & Pembaruan arsitektur, perubahan prinsip/framework, \emph{Request for
Architecture Work} baru (memulai siklus ADM berikutnya), dan \emph{Architecture
Contract} yang diperbarui. \\
Contoh (Restoran Go-Online) & \emph{Input:} Change Request (mitra GoFood minta
integrasi API; usulan ekspansi cabang); metrik order \& keluhan. \emph{Langkah:}
pantau $\rightarrow$ nilai dampak $\rightarrow$ klasifikasikan $\rightarrow$ tangani.
\emph{Output:} GoFood = \emph{Incremental} (adaptor baru); ekspansi multi-cabang =
\emph{Re-architecting} $\rightarrow$ \emph{Request for Architecture Work}. \\
Contoh (Bank Wanua) & \emph{Input:} Change Request (SNAP BI versi baru; usulan kanal
QRIS); metrik volume transaksi. \emph{Langkah:} pantau $\rightarrow$ nilai dampak
$\rightarrow$ klasifikasikan $\rightarrow$ tangani. \emph{Output:} SNAP BI baru =
\emph{Incremental} (adaptor diperbarui); QRIS = \emph{Re-architecting} $\rightarrow$
\emph{Request for Architecture Work}. \\
\hline
\end{tabularx}
\end{table}

\begin{table}[htbp]
\centering
\scriptsize
\caption{Input, langkah, dan output Requirements Management.}
\label{tab:io-rm}
\begin{tabularx}{\textwidth}{@{}p{0.17\textwidth} X@{}}
\hline
\textbf{Aspek} & \textbf{Uraian} \\
\hline
Input & Requirement dari \emph{semua} fase ADM \& \emph{Architecture Requirements
Specification}, ditambah requirement yang berubah beserta \emph{Requirements Impact
Assessment}. \\
Langkah & (1) Baseline \& dokumentasikan requirement; (2) simpan di
\emph{Architecture Requirements Repository}; (3) pantau \& tangkap requirement yang
berubah; (4) nilai dampaknya (\emph{Requirements Impact Assessment}); (5) perbarui
repository \& jaga \emph{traceability} lintas-fase. \\
Output & \emph{Requirements Impact Assessment}, \emph{Architecture Requirements
Specification} yang diperbarui, dan \emph{Architecture Requirements Repository} yang
terisi \& tertelusur. \\
Contoh (Restoran Go-Online) & \emph{Input:} requirement ``estimasi waktu antar
akurat'' \& ``data alergi tercatat''; Architecture Requirements Repository.
\emph{Langkah:} baseline $\rightarrow$ simpan $\rightarrow$ pantau perubahan
$\rightarrow$ \emph{Requirements Impact Assessment} saat menambah kanal.
\emph{Output:} matriks traceability (estimasi $\rightarrow$ mesin estimasi
$\rightarrow$ log akurasi). \\
Contoh (Bank Wanua) & \emph{Input:} requirement ``skrining AML wajib'' \& ``jejak
audit lengkap''; Architecture Requirements Repository. \emph{Langkah:} baseline
$\rightarrow$ simpan $\rightarrow$ pantau $\rightarrow$ \emph{Requirements Impact
Assessment} saat SNAP BI berubah. \emph{Output:} matriks traceability AML (kebijakan
$\rightarrow$ Mesin AML $\rightarrow$ WP2 $\rightarrow$ jejak audit). \\
\hline
\end{tabularx}
\end{table}

\section{Konteks Indonesia}

Di sektor keuangan Indonesia, governance dan change management bukan pilihan,
melainkan kewajiban. Penyelenggara sistem elektronik wajib menjaga keandalan,
keamanan, dan jejak sistem, serta melaporkan perubahan signifikan
\cite{pp712019}. Pembaruan standar pembayaran nasional seperti SNAP BI kerap
datang sebagai \emph{change driver} eksternal yang harus diserap tanpa mengganggu
layanan \cite{snapbi}. Perlindungan data pribadi menuntut bahwa setiap perubahan
yang menyentuh data nasabah dinilai dampaknya sejak awal, bukan setelah insiden
\cite{uupdp2022}. Tata kelola data lintas-instansi juga makin dituntut selaras
dengan prinsip berbagi-pakai yang terkelola \cite{satudata}. Karena itu,
\emph{Architecture Contract} dan matriks \emph{traceability} bukanlah formalitas;
keduanya justru menjadi bukti kepatuhan yang diminta OJK dan Bank Indonesia ketika
mengaudit rel transfer.

\section{Studi Kasus: Menjaga Arsitektur Tetap Patuh}

Governance, change management, dan traceability bekerja di domain apa pun. Bagian
ini menyajikan dua contoh: restoran yang \emph{go-online}---dari offline ke
online---sebagai contoh ringkas terlebih dahulu, lalu rel transfer Bank Wanua yang
lebih mendalam.

\subsection{Contoh Ringkas: Restoran Go-Online}

Sebuah restoran cepat saji yang selama ini hanya melayani makan di tempat
(\emph{offline}) kini ingin membuka kanal \emph{pemesanan online}. Transformasi
yang tampak kecil ini pun melewati ketiga proses berikut. \textbf{Fase G:} sebelum aplikasi diluncurkan, Compliance Review
menemukan modul pembayaran belum memenuhi standar keamanan kartu; Architecture
Board memberi \emph{Dispensation} satu bulan (sementara memakai payment gateway
pihak ketiga) disertai \emph{Change Request} integrasi pembayaran internal.
\textbf{Fase H:} permintaan integrasi ke mitra pengantaran (mis. GoFood) dinilai
\emph{Incremental}---cukup menambah adaptor dan memperbarui contract; sebaliknya,
ekspansi ke banyak cabang dengan dapur dan inventori terpusat adalah
\emph{Re-architecting} yang memicu siklus ADM baru. \textbf{Requirements
Management:} requirement ``estimasi waktu antar akurat'' ditelusuri dari sasaran
kepuasan pelanggan sampai mesin estimasi dan log akurasinya, sehingga ketika kanal
baru ditambahkan, dampaknya pada requirement lama langsung terlihat.

\subsection{Studi Kasus Utama: Rel Transfer Bank Wanua}

Setelah peta jalan Bab~\ref{ch:implementation-migration} dijalankan, rel transfer
Bank Wanua mulai beroperasi dan terus berubah. Tiga situasi berikut memperlihatkan
governance, change management, dan traceability bekerja bersama.

\textbf{Situasi 1 (Fase G).} Sebelum go-live SNAP BI, WP3 (Integrasi BI-FAST \&
Core) menjalani Compliance Review terhadap Architecture Contract. Adaptor BI-FAST
lulus, tetapi adaptor core banking masih menulis log tanpa jejak audit lengkap
karena keterbatasan mainframe. Architecture Board memberi \emph{Dispensation} tiga
bulan dengan syarat pemasangan lapisan pencatatan tambahan---bukan mengabaikan
penyimpangan, melainkan mengelolanya secara eksplisit.

\textbf{Situasi 2 (Fase H).} Enam bulan setelah go-live, Bank Indonesia merilis
versi baru SNAP BI. Tim menilai dampaknya: format pesan berubah, tetapi pola
integrasi tetap. Ini diklasifikasi sebagai perubahan \emph{Incremental}---contract
diperbarui, komponen adaptor disesuaikan, dan Compliance Review diulang, tanpa
merombak arsitektur. Sebaliknya, ketika manajemen mengusulkan menambah kanal
pembayaran QRIS dan \emph{virtual account}, penilaian menunjukkan kebutuhan node,
layanan, dan model data baru---sebuah \emph{Re-architecting} yang memicu siklus ADM
baru, kembali ke Fase A.

\textbf{Situasi 3 (Requirements Management).} Saat OJK meminta bukti bahwa kontrol
AML benar-benar diterapkan, tim membuka matriks traceability
(Gambar~\ref{fig:requirements-hub}): requirement ``skrining AML wajib'' tertaut ke
sasaran motivasi (Bab~\ref{ch:strategi-bisnis}), diwujudkan komponen Mesin AML \&
Limit (Bab~\ref{ch:application-architecture}), dibangun oleh WP2
(Bab~\ref{ch:implementation-migration}), dan menghasilkan jejak audit pada artefak
basis data (Bab~\ref{ch:technology-architecture}). Satu tarikan benang menautkan
kebijakan sampai bukti.

\begin{figure}[htbp]
\centering
\scalebox{0.9}{
\begin{tikzpicture}[
  font=\scriptsize,
  req/.style={rectangle, rounded corners=1.5mm, draw=orange!80!black, fill=orange!16,
    thick, align=center, text width=30mm, minimum height=9mm},
  lay/.style={rectangle, rounded corners=1.5mm, draw=gray!60, fill=gray!8,
    thick, align=center, text width=30mm, minimum height=9mm},
  link/.style={-Latex, semithick, draw=orange!70!black}
]
  \node[req] (r) {Requirement:\\``Skrining AML wajib''};
  \node[lay, above=6mm of r] (mot) {Sasaran motivasi\\(Bab~\ref{ch:strategi-bisnis})};
  \node[lay, right=10mm of r] (app) {Komponen Mesin AML\\(Bab~\ref{ch:application-architecture})};
  \node[lay, below=6mm of r] (wp) {WP2 Bangun Mesin AML\\(Bab~\ref{ch:implementation-migration})};
  \node[lay, left=10mm of r] (evd) {Jejak audit basis data\\(Bab~\ref{ch:technology-architecture})};
  \draw[link] (r) -- node[right,font=\tiny]{melatari} (mot);
  \draw[link] (r) -- node[above,font=\tiny]{diwujudkan} (app);
  \draw[link] (r) -- node[right,font=\tiny]{dibangun} (wp);
  \draw[link] (r) -- node[above,font=\tiny]{dibuktikan} (evd);
\end{tikzpicture}
}
\caption{\emph{Traceability} requirements sebagai pusat: satu requirement AML
tertaut ke motivasi, komponen aplikasi, work package, dan bukti audit di teknologi.}
\label{fig:requirements-hub}
\end{figure}

Diagram di atas menelusuri \emph{satu} requirement. Dalam praktik, keterlacakan
direkam sebagai \textbf{matriks traceability} yang memuat banyak requirement
sekaligus (Tabel~\ref{tab:traceability-matrix}): tiap baris menautkan sebuah
requirement ke sasaran, komponen/layanan, work package, dan bukti verifikasinya---
baik untuk rel transfer bank maupun restoran go-online.

\begin{table}[htbp]
\centering
\scriptsize
\caption{Contoh matriks traceability: tiap requirement ditelusuri lintas-lapisan.}
\label{tab:traceability-matrix}
\begin{tabularx}{\textwidth}{@{}p{0.19\textwidth} p{0.15\textwidth} p{0.16\textwidth} p{0.075\textwidth} X@{}}
\hline
\textbf{Requirement} & \textbf{Sasaran (motivasi)} & \textbf{Komponen/Layanan} &
\textbf{WP} & \textbf{Bukti/Verifikasi} \\
\hline
Skrining AML wajib \emph{(bank)} & Kepatuhan AML \& regulator & Mesin AML \& Limit &
WP2 & Jejak audit basis data \\
Transfer real-time \emph{(bank)} & Layanan cepat UMKM & Adaptor BI-FAST & WP3 & Log
latensi saat go-live \\
Estimasi antar akurat \emph{(restoran)} & Kepuasan pelanggan & Mesin Estimasi Antar &
Online & Log akurasi estimasi \\
Pembayaran online aman \emph{(restoran)} & Kepercayaan pelanggan & Modul Pembayaran &
Online & Sertifikasi keamanan kartu \\
\hline
\end{tabularx}
\end{table}

\section{Langkah Analisis dan Artefak Pendukung}

\subsection{Langkah 1: Susun Architecture Contract}
Rumuskan standar, prinsip, dan batasan yang mengikat tim implementasi---turunkan
dari prinsip arsitektur dan requirement kunci.

\subsection{Langkah 2: Tetapkan Gerbang Compliance Review}
Tentukan titik Compliance Review (mis. sebelum go-live) dan kriteria lulus/tidak,
termasuk mekanisme Dispensation.

\subsection{Langkah 3: Definisikan Protokol Perubahan}
Tetapkan bagaimana change driver ditangkap, dinilai dampaknya, dan diklasifikasi
menjadi \emph{Simplification}, \emph{Incremental}, atau \emph{Re-architecting}.

\subsection{Langkah 4: Bangun Matriks Traceability}
Tautkan tiap requirement kunci ke motivasi, komponen/layanan, work package, dan
bukti/verifikasinya.

\subsection{Langkah 5: Tentukan Metrik dan Value Realization}
Pilih indikator yang menunjukkan arsitektur memberi kapabilitas yang dijanjikan
(mis. waktu transfer, tingkat kepatuhan, insiden).

\subsection{Langkah 6: Petakan Governance dan Peran}
Tetapkan siapa yang memutuskan apa---Architecture Board, pemilik risiko, pemilik
produk---beserta irama peninjauannya.

\subsection{Langkah 7: Kaitkan dengan Regulasi}
Tautkan kontrol dan pelaporan perubahan dengan PP 71, UU PDP, dan standar SNAP BI.

\subsection{Langkah 8: Tulis Narasi Governance}
Ubah contract, protokol, dan matriks menjadi paragraf yang menjelaskan bagaimana
arsitektur dijaga tetap patuh dan relevan.

\begin{table}[htbp]
\centering
\scriptsize
\caption{Artefak pendukung Bab~10 dan fungsinya dalam makalah.}
\label{tab:artefak-bab9-gov}
\begin{tabularx}{\textwidth}{@{}p{0.26\textwidth} p{0.30\textwidth} X@{}}
\hline
\textbf{Artefak} & \textbf{Isi minimum} & \textbf{Dipakai untuk bagian makalah} \\
\hline
Architecture charter \& contract & Prinsip, standar, batasan, peran & Tata kelola
\& implikasi manajerial. \\
Protokol change management & Change driver, penilaian dampak, klasifikasi & Tata
kelola \& risiko. \\
Matriks traceability requirements & Requirement--arsitektur--WP--bukti & Analisis,
sintesis, keterlacakan end-to-end. \\
Kerangka metrik \& value realization & Indikator kapabilitas \& kepatuhan & Diskusi
\& implikasi. \\
\hline
\end{tabularx}
\end{table}

\section{Integrasi ke Makalah Publikasi}

Bab ini mengisi bagian \textbf{Tata Kelola}, \textbf{Risiko}, dan \textbf{Implikasi
Manajerial}, sekaligus menyediakan \emph{traceability} yang memperkuat seluruh
\textbf{Analisis}. Contoh paragraf:

\begin{quote}
Agar rel transfer tetap patuh dan relevan, arsitektur dikelola sebagai praktik
yang terus hidup. Sebuah Architecture Contract mengikat tiap work package pada
standar keamanan dan integrasi, dan Compliance Review menjadi gerbang sebelum
go-live SNAP BI. Perubahan eksternal---pembaruan SNAP BI, kanal pembayaran baru---
dinilai dan diklasifikasi: yang \emph{Incremental} diserap lewat governance, yang
mendasar memicu iterasi arsitektur baru. Matriks traceability memastikan setiap
kontrol, dari AML sampai perlindungan data, dapat ditelusuri dari kebijakan hingga
bukti audit, sehingga rekomendasi ini dapat dipertanggungjawabkan kepada OJK dan
Bank Indonesia.
\end{quote}

\section{Diskusi Kritis, Trade-off, dan Perangkap Umum}

\begin{table}[htbp]
\centering
\scriptsize
\caption{Trade-off dalam governance dan change management.}
\label{tab:tradeoff-bab9-gov}
\begin{tabularx}{\textwidth}{@{}p{0.24\textwidth} X X@{}}
\hline
\textbf{Keputusan} & \textbf{Jika terlalu ringan} & \textbf{Jika terlalu berat} \\
\hline
Ketatnya compliance & Arsitektur tergerus, utang teknis & Birokrasi lambat,
pengembang frustrasi. \\
Kebijakan Dispensation & Penyimpangan jadi permanen & Tak ada ruang gerak, proyek
mandek. \\
Ambang siklus ADM baru & Perubahan besar dipaksa \emph{Incremental} & Tiap
perubahan kecil memicu ADM, boros. \\
Kedalaman traceability & Tak teraudit, sulit dipertahankan & Overhead dokumentasi
berlebih. \\
\hline
\end{tabularx}
\end{table}

Perangkap umum: (1) governance yang hanya di atas kertas tanpa gerbang Compliance
Review nyata; (2) Dispensation tanpa batas waktu sehingga menjadi utang permanen;
(3) mencampur perubahan \emph{Incremental} dengan \emph{Re-architecting} sehingga
arsitektur berubah tanpa iterasi yang sah; dan (4) matriks traceability yang tidak
dirawat sehingga usang saat audit tiba.

\section{Latihan Terapan Individual}

\subsection*{Latihan~10.1: Architecture Contract dan Gerbang Compliance}
Untuk organisasi studi kasus Anda, susun Architecture Contract ringkas (prinsip,
standar, batasan) dan tetapkan minimal satu gerbang Compliance Review beserta
kriterianya.

\subsection*{Latihan~10.2: Protokol Change Management}
Rumuskan cara menangkap change driver dan klasifikasikan tiga contoh perubahan
(\emph{Simplification}, \emph{Incremental}, \emph{Re-architecting}) beserta
keputusannya.

\subsection*{Latihan~10.3: Matriks Traceability}
Ambil dua requirement kunci dan telusuri masing-masing ke motivasi, komponen,
work package, dan bukti verifikasinya.

\subsection*{Latihan~10.4: Paragraf Governance}
Tulis satu paragraf yang menjelaskan bagaimana arsitektur Anda dijaga tetap patuh
dan relevan, termasuk peran, metrik, dan kaitan regulasi.

\begin{table}[htbp]
\centering
\scriptsize
\caption{Rubrik latihan Bab~10.}
\label{tab:rubrik-bab9-gov}
\begin{tabularx}{\textwidth}{@{}p{0.25\textwidth} X X@{}}
\hline
\textbf{Komponen} & \textbf{Indikator baik} & \textbf{Perlu diperbaiki bila} \\
\hline
Contract \& compliance & Kriteria konkret \& gerbang jelas & Prinsip normatif tanpa
gerbang atau kriteria. \\
Protokol perubahan & Klasifikasi \& keputusan beralasan & Perubahan tanpa penilaian
dampak atau kelas. \\
Traceability & Requirement tertaut penuh \& berbukti & Tautan terputus atau tanpa
verifikasi. \\
Narasi governance & Menjelaskan peran, metrik, \& regulasi & Deskriptif tanpa
akuntabilitas. \\
\hline
\end{tabularx}
\end{table}

\section{Ringkasan, Refleksi, dan Penutup Metode}

\subsection*{Poin Kunci}
\begin{itemize}
  \item \textbf{Fase G (Implementation Governance)} menjaga kepatuhan implementasi
        lewat Architecture Contract, Compliance Review, dan Dispensation berbatas
        waktu.
  \item \textbf{Fase H (Architecture Change Management)} memantau change driver dan
        mengklasifikasinya: \emph{Simplification}, \emph{Incremental}, atau
        \emph{Re-architecting} (memicu ADM baru).
  \item \textbf{Requirements Management} adalah pusat ADM yang menjaga
        \emph{traceability} dari motivasi sampai bukti solusi.
  \item Bersama-sama ketiganya menjadikan arsitektur \emph{hidup, terkelola, dan
        dapat diaudit}, bukan dokumen sekali jadi.
  \item Dalam konteks Indonesia, contract dan traceability adalah bukti kepatuhan
        yang diminta OJK, Bank Indonesia, dan regulasi PDP.
\end{itemize}

\subsection*{Uji Mandiri}
\begin{enumerate}
  \item Jelaskan perbedaan Compliance Review (Fase G) dan Architecture Change
        Management (Fase H).
  \item Kapan sebuah perubahan cukup ditangani lewat governance, dan kapan ia
        harus memicu siklus ADM baru?
  \item Ambil satu requirement pada kasus Anda dan telusuri tautannya dari
        motivasi sampai bukti solusi.
\end{enumerate}

\subsection*{Penutup: Dari Rancangan ke Praktik yang Hidup}
Ketujuh lapisan sebelumnya membangun arsitektur enterprise yang utuh: dari motivasi
turun ke strategi, bisnis, data, aplikasi, teknologi, lalu peta jalan implementasi.
Bab ini menutup \emph{metode}-nya---menunjukkan bahwa arsitektur yang baik bukan
sekadar rancangan yang indah, melainkan praktik yang \emph{dijaga} dan
\emph{ditelusuri}. Governance memastikan yang dibangun setia pada rancangan;
change management memastikan arsitektur menyesuaikan diri dengan dunia yang
bergerak; dan Requirements Management menjaga agar setiap keputusan selalu dapat
ditarik kembali ke alasannya. Dengan menutup roda ADM, mahasiswa tidak hanya
merancang satu arsitektur target, tetapi juga merancang \emph{cara merawatnya}---
kompetensi yang membedakan arsitek enterprise dari sekadar pembuat diagram.

\section{Bacaan Lanjutan dan Lampiran}

\subsection*{Bacaan Inti}
\begin{itemize}
  \item The Open Group untuk Fase G, H, dan Requirements Management ADM
        \cite{togaf10}.
  \item Ross, Weill, \& Robertson untuk \emph{IT governance} dan fondasi eksekusi
        \cite{ross2006enterprise,ross2019designed}.
\end{itemize}

\subsection*{Bacaan untuk Memperdalam}
\begin{itemize}
  \item Newman untuk umpan balik, \emph{continuous delivery}, dan perubahan yang
        dikelola \cite{newman2021building}.
  \item Untuk konteks Indonesia: penyelenggaraan sistem elektronik \cite{pp712019},
        perlindungan data pribadi \cite{uupdp2022}, standar API pembayaran
        \cite{snapbi}, dan tata kelola data \cite{satudata}.
\end{itemize}

\subsection*{Lampiran Kerja Bab~10}
\begin{description}
  \item[Architecture charter \& contract] Prinsip, standar, batasan, peran.
  \item[Protokol change management] Change driver, penilaian dampak, klasifikasi.
  \item[Matriks traceability] Requirement--arsitektur--WP--bukti.
  \item[Kerangka metrik] Indikator kapabilitas, kepatuhan, dan value realization.
\end{description}
\setlength{\emergencystretch}{0pt}
